\documentclass[12pt,aspectratio=169]{beamer}
\input{header_beamer}
\usetheme{Boadilla}
\usecolortheme{beaver}
\setbeamertemplate{page number in head/foot}{}
\setbeamertemplate{navigation symbols}{}

\title{Lecture-03: Independence}
\date{}

\begin{document}
\pagestyle{empty}
\maketitle

\begin{frame}{Law of Total Probability}
\begin{block}{Exercise : Countably infinite coin tosses} 
\begin{itemize}
    \item <1-> Consider a sequence of coin tosses, such that the sample space is $\Omega = \set{H,T}^\N$. 
    \item <2-> For set of outcomes $E_n \triangleq \set{\omega \in \Omega: \omega_n = H}$, consider an event space generated by $\sF \triangleq \sigma(\set{E_n: n\in \N})$. 
    \item <3-> $\sF_n \triangleq \sigma(\set{E_i: i \in [n]})$
    \item <4-> $A_n \triangleq \set{\omega \in \Omega: \omega_i = H \text{ for some } i \in [n]} = \cup_{i=1}^nE_i \in \sF$ 
    \item <5-> $B_n \triangleq \set{\omega \in \Omega: \omega_1 = \dots = \omega_{n-1}= T, \omega = H} = \cap_{i=1}^{n-1}E_i^c\cap E_n \in \sF$
    \begin{enumerate}
        \item <6-> Show that $\sF = \sigma(\set{\sF_n: n \in \N})$. 
        \item <7-> Show that $\sigma(\set{A_n: n \in \N}) \subseteq \sF$ and $\sigma(\set{B_n: n\in \N}) \subseteq \sF$. 
    \end{enumerate}
\end{itemize}
\end{block}
\end{frame}

\begin{frame}{Law of Total Probability}
\begin{block}{Theorem : Law of total probability} 
For a probability space $(\Omega, \sF, P)$, 
consider a sequence of events $B\in \sF^\N$ that partitions the sample space $\Omega$, i.e. 
$B_m \cap B_n = \emptyset$ for all $m \neq n$, and $\cup_{n \in \N}B_n = \Omega$. 
Then, for any event $A \in \sF$, we have 
\EQ{
P(A) = \sum_{n \in \N}P(A \cap B_n).
}
\end{block}
\end{frame}
\begin{frame}{Law of Total Probability}
\begin{Proof}
\begin{itemize}
    \item <1-> \EQ{
                    A = A \cap \Omega = A \cap (\cup_{n \in \N}B_n) = \cup_{n \in \N}(A \cap B_n).
                    }
    \item <2-> The sequence $(A \cap B_n \in \sF: n \in \N)$ is mutually disjoint.
    \item <3-> The result follows from the countable additivity of probability of disjoint events.
\end{itemize}
\end{Proof}
\end{frame}
\begin{frame}{Law of Total Probability}
\begin{Example}[Countably infinite coin tosses] 
Consider the sample space $\Omega = \set{H,T}^\N$ and event space $\sF$ generated by sequence $E \in \sF^\N$. 
We observe that any event $A \in \sF_n$ can be written as 
\begin{equation*}
A = \cup_{\omega \in A}\set{\omega} 
= \cup_{\omega \in A}\cap_{i=1}^n(\set{\omega \in E_i}\cup\set{\omega \notin E_i}).
\end{equation*}  
\end{Example}
\end{frame}



\begin{frame}{Independence}
\begin{Definition}[Independence of events]<1->
For a probability space $(\Omega, \sF, P)$, a family of events $A \in \sF^I$ is said to be independent, 
if for any finite set $F \subseteq I$, we have 
\EQ{
P(\cap_{i \in F}A_i) = \prod_{i \in F}P(A_i). 
}
\end{Definition}
\begin{block}{Reminder}<2->
The certain event $\Omega$ and the impossible event $\emptyset$ are always independent to every event $A \in \sF$. 
\end{block}
\end{frame}
\begin{frame}


\begin{Example}[Two coin tosses] 
\begin{itemize}
\item <1-> Consider two coin tosses with sample space $\Omega = \set{HH, HT, TH, TT}$, and event space $\sF = \cP(\Omega)$. 
\item <2-> Define $P$
\EQ{
P(\set{HH}) = P(\set{HT}) = P(\set{TH}) =P(\set{TT}) = \frac{1}{4}. 
}
\item <2-> Let event $E_1 \triangleq \set{HH, HT}$ and $E_2 \triangleq \set{HH, TH}$ correspond to getting a head on the first or the second toss respectively. 
\item <3-> $P(E_1) = P(E_2) = \frac{1}{2}$ 
\item <4-> The intersecting event $E_1\cap E_2 = \set{HH}$ with the probability $P(E_1 \cap E_2) = \frac{1}{4}$. 
\item <5-> 
\EQ{
P(E_1 \cap E_2) = P(E_1)P(E_2).
}
That is, events $E_1$ and $E_2$ are independent. 
\end{itemize}
\end{Example}
\end{frame}

\begin{frame}
\begin{Example}[Countably infinite coin tosses] 
\begin{itemize}
\item <1-> Consider the outcome space $\Omega = \set{H,T}^\N$ and event space $\sF$ generated by the sequence $E$.  
\item <2-> Define $P:\sF \to [0,1]$ as $P(\cap_{i\in F}E_i)= p^{\abs{F}}$ for any finite subset $F \subseteq \N$.  
\item <3-> $E \in \sF^\N$ is a sequence of independent events. 
\item <4-> Consider $A, B \in \sF^\N$, 
where $A_n \triangleq \cup_{i=1}^nE_i$ and $B_n \triangleq \cap_{i=1}^{n-1}E_i^c\cap E_n \in \sF$ for all $n\in\N$.
\item <5-> $P(A_n) = 1-(1-p)^n$ and $P(B_n) = p(1-p)^{n-1}$ for $n \in \N$. 
\item <6-> For any $\omega \in \Omega$, define $k_n(\omega) \triangleq\sum_{i=1}^n\SetIn{H}(\omega_i) = \sum_{i=1}^n\SetIn{\omega \in E_i}$. 
\item <7-> For $A \in \sF_n = \sigma(\set{E_i: i \in [n]})$,  
\EQ{
P(A) = \sum_{\omega \in A} \prod_{i=1}^n\Big[P\set{\omega \in E_i}  + P\set{\omega \in E_i^c}\Big] 
= \sum_{\omega \in A}p^{k_n(\omega)}(1-p)^{n-k_n(\omega)}.
}  
\end{itemize}
\end{Example}
\end{frame}

\begin{frame}{Independence}
\begin{Example}[Counter example] 
\begin{itemize}
\item <1-> Consider a probability space $(\Omega, \sF, P)$ and the events $A_1, A_2, A_3 \in \sF$. 
\item <2-> $P(A_1\cap A_2\cap A_3) = P(A_1)P(A_2)P(A_3)$ is not sufficient to guarantee the independence of the three events. 
\item <3-> In particular, if
\meq{2}{
&P(A_1\cap A_2\cap A_3) = P(A_1)P(A_2)P(A_3),&&P(A_1\cap A_2\cap A_3^c) \neq P(A_1)P(A_2)P(A_3^c),
}\\
\item <4-> Then,\\ $P(A_1\cap A_2) =  P(A_1\cap A_2\cap A_3) + P(A_1\cap A_2\cap A_3^c) \neq P(A_1)P(A_2)$.
\end{itemize}
\end{Example} 
\end{frame}

\begin{frame}{Independence for a family of collections of events}
\begin{Definition}
A family of collections of events $(\sA_i \subseteq \sF: i \in I)$ is called independent,
if for any finite set $F \subseteq I$ and $A_i \in \sA_i$ for all $i \in F$, we have 
\EQ{
P(\cap_{i \in F}A_i) = \prod_{i \in F}P(A_i).
} 
\end{Definition}
\end{frame}
\begin{frame}{Conditional Probability}
\begin{itemize}
\item <1-> Consider $N$ trials of a random experiment over an outcome space $\Omega$ and an event space $\sF$. $\omega_n \in \Omega$ denotes the outcome of the experiment of the $n$th trial. 
\item <2-> $A, B \in \sF$
\item <3-> $N(A)$, $N(B)$ and $N(A\cap B)$. 
\end{itemize}
\begin{block}{Equations}<4->
\meq{3}{
&N(A) = \sum_{n=1}^N\SetIn{\omega_n \in A}, &&N(B) = \sum_{n=1}^N\SetIn{\omega_n \in B}, &&N(A\cap B) = \sum_{n=1}^N\SetIn{\omega_n \in A \cap B}.
}
\end{block}
\end{frame}
\begin{frame}{Conditional Probability}
\begin{itemize}
\item <1-> Relative frequency of events $A, B, A\cap B$ in $N$ trials are $\frac{N(A)}{N}, \frac{N(B)}{N}, \frac{N(A\cap B)}{N}$ respectively.  
\item <2-> We can find the relative frequency of events $A$, on the trials where $B$ occurred as 
\begin{block}{Relative frequency}
\EQ{
\frac{\frac{N(A\cap B)}{N}}{\frac{N(B)}{N}} = \frac{N(A \cap B)}{N(B)}.
} 
\end{block}
\end{itemize}
\end{frame}
\begin{frame}{Conditional Probability}
\begin{Definition}[Conditional Probability]<1->
Fix an event $B \in \sF$ such that $P(B) > 0$, we can define the conditional probability $P(\cdot | B): \sF \to [0,1]$ of any event $A \in \sF$ conditioned on the event $B$ as 
\EQ{
P(A|B) = \frac{P(A\cap B)}{P(B)}.
}
\end{Definition}
\begin{block}{Lemma:Conditional probability}<2->
For any event $B \in \sF$ such that $P(B) > 0$, the conditional probability $P(\cdot | B): \sF \to [0,1]$ is a probability measure on space $(\Omega, \sF)$. 
\end{block}
\end{frame}
\begin{frame}{Conditional Probability}
\begin{Proof} 
\begin{enumerate}
\item <1-> [\textbf{Non-negativity:}] For all events $A \in \sF$, we have $P(A|B) \ge 0$ since $P(A \cap B) \ge 0$. 
\item <2-> [\textbf{$\sigma$-additivity:}] For an infinite sequence of mutually disjoint events $(A_i \in \sF: i \in \N)$ such that $A_i \cap A_j = \emptyset$ for all $i \neq j$, we have $P(\cup_{i \in \N}A_i|B) = \sum_{i \in \N}P(A_i|B)$. 
This follows from disjointness of the sequence $(A_i\cap B \in \sF: i \in \N)$. 
\item<3->  [\textbf{Certainty:}] Since $\Omega \cap B = B$, we have $P(\Omega |B) = 1$. 
\end{enumerate} 
\end{Proof}
\end{frame}
\begin{frame}{Conditional Probability}
\begin{block}{Reminder}
\begin{itemize}
\item <1-> For two independent events $A, B \in \sF$ such that $P(A\cap B) > 0$, 
we have  $P(A|B) = P(A)$ and  $P(B|A) = P(B)$. 
If either $P(A) = 0$ or $P(B) = 0$, then $P(A \cap B) = 0$. 
\item <2-> For any partition $B$ of the sample space $\Omega$, if $P(B_n) > 0$ for all $n \in \N$, then from the law of total probability and the definition of conditional probability, we have 
\EQ{
P(A) = \sum_{n \in \N}P(A|B_n)P(B_n).
}
\end{itemize}
\end{block}
\end{frame}



\begin{frame}{Conditional Independence}
\begin{Definition}[Conditional independence of events]
For a probability space $(\Omega, \sF, P)$, a family of events $A \in \sF^I$ is said to be conditionally independent given an event $C \in \sF$ such that $P(C) > 0$, 
if for any finite set $F \subseteq I$, we have 
\EQ{
P(\cap_{i \in F}A_i|C) = \prod_{i \in F}P(A_i|C). 
}
\end{Definition}
\end{frame}
\begin{frame}{Conditional Independence}
\begin{block}{Reminder}
\begin{itemize}

\item <1-> Let $C \in \sF$ be an event such that $P(C) > 0$. 
Two events $A, B \in \sF$ are said to be conditionally independent given event $C$, if 
\EQ{
P(A \cap B|C) = P(A|C) P(B|C). 
}
If the event $C = \Omega$, it implies that $A, B$ are independent events. 

\item <2-> Two events may be independent, but not conditionally independent and vice versa. 
\end{itemize}
\end{block}
\end{frame}
\begin{frame}
\begin{Example}[Conditional Independence]
\begin{itemize}
\item <1-> Two independent events $A, B \in \sF$ such that $P(A\cap B) > 0$ and $P(A \cup B) < 1$. 
\item <2-> $A$ and $B$ are not conditionally independent given $A\cup B$.  
\EQ{
P(A \cap B | A\cup B) = \frac{P(A \cap B)}{P(A\cup B)} 
= \frac{P(A)P(B)}{P(A \cup B)} =  {P(A|A \cup B)}{P(B)}. 
}
\item <3-> $P(B|A\cup B) = \frac{P(B)}{P(A\cup B)} \neq P(B)$ \implies $P(A \cap B | A\cup B) \neq P(A|A\cup B)P(B|A\cup B)$. 
\item <4-> Two non-independent events $A, B \in \sF$ such that $P(A) > 0$. 
\item <5-> $A$ and $B$ are conditionally independent given $A$. 
\EQ{
P(A \cap B | A) = \frac{P(A \cap B)}{P(A)} = P(B|A)P(A|A). 
}
\end{itemize}
\end{Example}
\end{frame}
\end{document}
